\phantomsection\addcontentsline{toc}{section}{Water Treatment}
\ECchapter{14}{Water Treatment}{How can engineering provide safe drinking water?}{ECBlue}

\ECObjectives{
\begin{itemize}
\item Identify the principal barriers in a drinking-water treatment chain.
\item Explain filtration, disinfection and reverse osmosis using physical principles.
\item Select and monitor a robust treatment solution from water-quality and local constraints.
\end{itemize}}

\begin{multicols}{2}
\begin{engineeringnews}
Membrane filtration, ultraviolet disinfection and lower-energy desalination are improving access to
safe water. Yet the best solution is not always the most advanced one: reliability, operator skills,
spare parts, energy availability and safe sludge or brine management often determine whether a plant
succeeds over its complete lifetime.
\end{engineeringnews}

\begin{ecarticle}
Raw water may contain suspended particles, microorganisms, organic matter, dissolved salts and
chemical pollutants. No single process removes every contaminant, so drinking-water plants combine
several barriers. The treatment chain is selected according to source-water quality, local
constraints and the required quality of the final water.

Screens first remove leaves and large debris. During coagulation, a chemical reagent destabilises
very small particles that would otherwise remain in suspension. Gentle mixing then creates larger
flocs. In a sedimentation tank, these flocs settle under gravity and form sludge that must be
collected, dewatered and treated.

Filtration removes the particles that remain. Sand filters trap material inside a porous bed,
whereas activated carbon also adsorbs some organic molecules responsible for taste, odour or
pollution. Membrane systems use pores of controlled size. As a filter becomes dirty, pressure loss
increases, so periodic backwashing or chemical cleaning is necessary.

Disinfection inactivates harmful microorganisms. Chlorine can maintain protection throughout the
distribution network, whereas ultraviolet light acts rapidly but leaves no disinfectant residual.
Engineers often combine methods and monitor turbidity, pH, conductivity, flow and disinfectant
concentration continuously.

Reverse osmosis is used for desalination. Pumps apply a pressure greater than the natural osmotic
pressure, forcing water through a semi-permeable membrane while most salts remain in a concentrated
brine. The process produces high-quality water but requires energy, pretreatment and careful brine
management.

For a remote community, maintainability may be more important than maximum theoretical performance.
Gravity flow, modular filters, solar pumping and simple sensors can reduce operating costs. A robust
system must also include safe storage and prevent treated water from being contaminated again.
\end{ecarticle}

\begin{tcolorbox}[title=\sffamily\bfseries Engineering Insight,colback=yellow!10,colframe=ECOrange]
Safe water depends on verified barriers, not on appearance. A clear sample may still contain
pathogens, while a highly sophisticated plant may fail if reagents, spare parts or trained operators
are unavailable. Treatment performance and maintainability must therefore be designed together.
\end{tcolorbox}

\begin{engineeringfocus}
\textbf{Multi-barrier treatment chain}
\begin{center}
\resizebox{\linewidth}{!}{%
\begin{tikzpicture}[font=\scriptsize,>=Latex,node distance=3.5mm]
\node[draw=ECBlue,fill=ECLightBlue,rounded corners,align=center] (raw) {raw\\water};
\node[draw=ECGreen,fill=ECGreen!10,rounded corners,align=center,right=of raw] (coag) {coagulation\\and flocculation};
\node[draw=ECOrange,fill=ECOrange!10,rounded corners,align=center,right=of coag] (sed) {sedimentation};
\node[draw=ECPurple,fill=ECPurple!10,rounded corners,align=center,below=of sed] (fil) {filtration};
\node[draw=ECRed,fill=ECRed!8,rounded corners,align=center,left=of fil] (dis) {disinfection};
\node[draw=ECBlue,fill=ECLightBlue,rounded corners,align=center,left=of dis] (tank) {safe storage\\and distribution};
\draw[->,thick] (raw)--(coag); \draw[->,thick] (coag)--(sed);
\draw[->,thick] (sed)--(fil); \draw[->,thick] (fil)--(dis); \draw[->,thick] (dis)--(tank);
\draw[->,thick,dashed,ECRed] (sed.south) -- +(0,-.65) node[below]{sludge treatment};
\end{tikzpicture}%
}
\end{center}
Each stage is a barrier. If one process performs less effectively, the remaining barriers reduce the
risk of unsafe water reaching consumers.
\end{engineeringfocus}

\begin{keyfigures}
\begin{tabularx}{\linewidth}{@{}Xr@{}}
Raw-water quality & highly variable\\
Filtered-water target & very low turbidity\\
Reverse-osmosis driving force & high pressure\\
Main membrane issue & fouling\\
Essential final barrier & safe storage
\end{tabularx}
\end{keyfigures}

\begin{tcolorbox}[title=\sffamily\bfseries Turbidity through the treatment chain,colback=white,colframe=ECBlue]
\begin{center}
\begin{tikzpicture}
\begin{axis}[ybar,width=.96\linewidth,height=45mm,ylabel={Relative turbidity (a.u.)},ymin=0,ymax=20,symbolic x coords={Raw water,Coagulation,Sedimentation,Filtration,Final water},xtick=data,x tick label style={rotate=38,anchor=east,font=\scriptsize},yticklabel style={font=\scriptsize},grid=major]
\addplot table[x=stage,y=turbidity,col sep=comma]{data/EC14/treatment_performance.csv};
\end{axis}
\end{tikzpicture}
\end{center}
\footnotesize Successive barriers progressively reduce suspended matter. The values are illustrative;
actual performance depends on source water and operating conditions.
\end{tcolorbox}

\begin{tcolorbox}[title=\sffamily\bfseries Reverse osmosis principle,colback=white,colframe=ECPurple]
\begin{center}
\begin{tikzpicture}[font=\scriptsize,>=Latex]
\draw[thick,fill=ECBlue!8] (0,0) rectangle (2.2,1.7);
\draw[thick,fill=ECGreen!8] (3.1,0) rectangle (5.3,1.7);
\draw[very thick,ECPurple] (2.65,-.15)--(2.65,1.85);
\node[rotate=90,ECPurple] at (2.45,.85){membrane};
\foreach \x/\y in {.4/.35,.8/.9,1.3/.55,1.7/1.25,.55/1.4,1.9/.45}{\node at (\x,\y){$+$};}
\draw[->,very thick,ECOrange] (-.4,.85)--(.15,.85) node[midway,above]{pressure};
\draw[->,very thick,ECBlue] (2.15,.85)--(3.05,.85);
\node[align=center] at (1.1,-.35){salty feed};
\node[align=center] at (4.2,-.35){freshwater};
\node[align=center,ECRed] at (1.1,2.05){concentrated brine remains};
\end{tikzpicture}
\end{center}
\end{tcolorbox}

\begin{vocabulary}
\begin{tabularx}{\linewidth}{@{}lX@{}}
raw water & eau brute\\
coagulation & coagulation\\
floc & floc\\
sedimentation & décantation\\
activated carbon & charbon actif\\
backwashing & lavage à contre-courant\\
disinfection & désinfection\\
turbidity & turbidité\\
reverse osmosis & osmose inverse\\
brine & saumure
\end{tabularx}
\end{vocabulary}

\begin{questions}
\item Why does a plant use several treatment barriers?
\item Explain the difference between coagulation and sedimentation.
\item What additional function can activated carbon provide?
\item Compare chlorine and ultraviolet disinfection.
\item Use the graph to identify the stage with the largest apparent turbidity reduction.
\item Why are pretreatment and maintenance important in reverse osmosis?
\item Which measurements should trigger an automatic alarm or plant shutdown?
\end{questions}

\begin{applicationexercise}
A small plant treats 500 m$^3$ of water per day. Raw water has a turbidity of 18 NTU. After
sedimentation, turbidity is reduced by 75\%, and the filter then removes 92\% of the remaining
turbidity.
\begin{enumerate}
\item Calculate the turbidity after sedimentation and filtration.
\item A target of 0.5 NTU is required. State whether the simplified chain meets it.
\item Calculate the daily volume that must be reprocessed if 3\% of production is rejected.
\item Explain why turbidity alone cannot prove microbiological safety.
\end{enumerate}
\end{applicationexercise}

\begin{engineeringchallenge}
Design a compact treatment system for a village beside a turbid river. Characterise the hazards,
choose the treatment barriers, energy source and monitoring sensors, and define acceptance limits.
Describe routine maintenance, reagent and spare-part needs, sludge management, safe storage and the
action taken when one measurement exceeds its limit. Compare a simple robust option with a more
automated membrane-based option and justify your recommendation.
\end{engineeringchallenge}

\begin{speaking}
Compare a central treatment plant with small local treatment units. Present a four-minute technical
recommendation based on water quality, energy, maintenance, resilience, cost and operator skills.
\end{speaking}
\end{multicols}

\subsection*{Sources}
\ECsource{World Health Organization -- Drinking Water}{https://www.who.int/news-room/fact-sheets/detail/drinking-water}
\ECsource{U.S. Environmental Protection Agency -- Drinking Water}{https://www.epa.gov/ground-water-and-drinking-water}
